\section{引言}
\label{sec:introduction}
土地覆盖模型常在一种制图产品上训练，再在另一种产品上评估。影像、传感器、分辨率与地理区域变化可能影响性能，而不同产品的图例差异也会改变模型实际获得的正例监督\citep{tuia2016domain,bruzzone2009domain,yang2017harmonization}。如果某个目标类别在源数据的精确类别标签中缺失，仅凭接近零的目标 IoU 很难解释迁移失败；这不足以证明影像域偏移是唯一原因。

本文将某次映射决策中最细的软件输出类别称为\emph{操作性叶类别}，并不据此声称已经建立完整一致的植物学分类体系。在明确的类别映射下，目标叶类别可能获得精确源支持、仅获得较粗祖先类别支持，或没有映射支持。这些是监督记录，而非彼此排斥的迁移误差原因。精确标签缺失也不意味着预训练模型不具备该概念的表征，或异质监督无法支持该概念；它只表示指定的源训练掩膜没有提供这一精确叶类别的正例，尽管损失函数仍会影响其预测。

标签空间协调与标注冲突已有研究基础：MSeg 协调类别体系和像素标注，UniSeg 处理多数据集标签偏移\citep{lambert2020mseg,kim2022uniseg}。分别报告共同与私有类别性能，以及考察负迁移，也有既有先例\citep{choe2024bus,choe2025unida,ren2026sats}。本文的问题更具体：原始掩膜支持审计能怎样改变对某次迁移失败的解释？在固定更新预算和源暴露控制下，所考察的源混合带来多大的恢复收益，又保留了多少源支持类别性能？

在冻结的编码表中，LoveDA 没有被分配到草本输出索引的原始类别，FLAIR 原始 ID 10 则提供该训练分配。OpenEarthMap 的 rangeland 是这一索引所采用的较宽目标代理类别，而非纯草本参考类别（补充材料 S2）。主要 CLIP 实验比较 LoveDA 单源基线与 LoveDA--FLAIR 混合，并保持目标评估总体一致。暴露控制与重复裁块控制约束了 LoveDA 采样量和更新表示的变化。历史原生 U-Net 对比在同一目标上检验另一种记录模型家族中的恢复与保持方向。已记录的错配影像—掩膜配置则提供单独的 CLIP 敏感性分析，其来源证据有限。辅助影像与标签同时进入训练，因此这些源混合对比均不能单独识别标签的因果作用。

本文的贡献是在上述限定条件下开展经验审计与诊断设计：将版本化的操作性类别映射关联到原始掩膜正例支持；在指定控制下同时报告缺失类别恢复与源支持类别损失；并在明确各自依赖关系及来源边界的前提下，考察历史模型家族扩展和反向源混合。主要对比中，rangeland IoU 的差值为+0.3118，支持类别 mIoU 的差值为-0.1011；历史 U-Net 记录重复了这一方向。校准适配器和 U-Net 是测量工具，而非本文提出的新算法。本文不主张首次发现标签不匹配或负迁移，也没有声称已验证解决该权衡的方法。
